\begin{teorema}[Regularidad a actividad fija]\label{thm:lip_interior}
Sean $\rho\in(0,1)$, $x=1-\rho$ y $\lambda=\lambda(\rho)$. Las versiones continuas de las vacancias,
las proyecciones y el maximizador cumplen
\[
 \Lip(u_f),\ \Lip(m_f)\le L(\lambda),\qquad \|\nu_\rho\|_\infty\le\lambda,\qquad
 \Lip_\infty(\nu_\rho)\le6\lambda L(\lambda),
\]
y, para $\Theta_x(r,c)=r_R(r)\,r_C(c)=u_R(r)u_C(c)/x^2$,
\[
 \|\Theta_x\|_\infty\le x^{-2},\qquad \Lip_\infty(\Theta_x)\le2L(\lambda)/x^2 ,
\]
donde $\Lip_\infty$ se refiere a la norma del máximo en el cuadrado. Las cotas incluyen todo el borde
y son uniformes en cada compacto de $(0,1)$.
\end{teorema}

\begin{proof}
La estimación~\eqref{eq:traslacion} da la primera cota para cualesquiera dos puntos del mismo intervalo.
Telescopar los cuatro factores de la ecuación de Euler cuesta
$\lambda L(\lambda)(|\Delta r|+|\Delta c|+|\Delta r-\Delta c|+|\Delta r+\Delta c|)\le6\lambda L(\lambda)\|\Delta(r,c)\|_\infty$,
y el producto de los dos factores de $\Theta_x$ cuesta $2L(\lambda)\|\Delta(r,c)\|_\infty/x^2$. En un
compacto de masas basta sustituir $\lambda$ por su máximo y $x$ por su mínimo, porque las cotas son
monótonas y $\lambda$ es creciente.
\end{proof}

El Teorema~\ref{thm:lip_interior} es la parte~(a) del Teorema~\ref{thm:lipschitz_intro}.

\begin{observacion}[Degeneración en los extremos]\label{obs:degeneracion}
Como $\lambda=b/x^2\le b_+/x^2$ para $x\le\frac12$ (Proposición~\ref{prop:cotas}), las constantes del
Teorema~\ref{thm:lip_interior} son $O(x^{-6})$ para las vacancias y $O(x^{-8})$ para $\nu$ y $\Theta_x$
cuando $x\downarrow0$; cuando $\rho\downarrow0$, $\lambda\le\rho(1+\lambda)^4=O(\rho)$ (Proposición~\ref{prop:concava}) da $O(\rho)$, $O(\rho^2)$ y $O(\rho)$,
respectivamente. Son cotas superiores, no exponentes óptimos. No se afirma regularidad $C^1$ global:
los dominios diagonales cambian de geometría y una cota Lipschitz no excluye posibles puntos angulosos.
\end{observacion}

\begin{proposicion}[Dependencia de la actividad]\label{prop:actividad}
Sean $0<\lambda,\lambda'\le\Lambda$ y $v_\lambda,v_{\lambda'}$ los campos de los maximizadores
correspondientes. Entonces
\begin{gather*}
 \|v_\lambda-v_{\lambda'}\|_\cH\le\sqrt6\,(1+\Lambda)|\lambda-\lambda'|,\qquad
 \|u_{f,\lambda}-u_{f,\lambda'}\|_\infty\le C_\Lambda|\lambda-\lambda'|,\\
 \|\nu_\lambda-\nu_{\lambda'}\|_\infty\le(1+4\Lambda C_\Lambda)|\lambda-\lambda'|,
\end{gather*}
con $C_\Lambda=1+5\Lambda(1+\Lambda)$.
\end{proposicion}

\begin{proof}
Las vacancias de ambos están en $[(1+\Lambda)^{-1},1]$. Como $\cG_{\lambda'}(v_{\lambda'})=0$,
$\cG_{\lambda,f}(v_{\lambda'})=(\lambda-\lambda')\,u_{f,\lambda'}J_f(u_{\lambda'})$, cuya norma en $\cH$
es a lo sumo $\sqrt6|\lambda-\lambda'|$, porque las longitudes de los cuatro dominios suman seis. El
Lema~\ref{lem:estabilidad} con $\kappa'=(1+\Lambda)^{-1}$ da la primera cota. Comparando
$u_{f,\lambda}=1/(1+\lambda J_f(u_\lambda))$ con su versión en $\lambda'$ y usando~\eqref{eq:conversion},
$\|u_{f,\lambda}-u_{f,\lambda'}\|_\infty\le|\lambda-\lambda'|+\Lambda\sqrt3\,\|v_\lambda-v_{\lambda'}\|_\cH\le C_\Lambda|\lambda-\lambda'|$.
La ecuación de Euler y el producto de cuatro factores dan la última cota.
\end{proof}

En particular, como $\lambda$ es continua, $x\mapsto(\nu_{1-x},u_{R,x},u_{C,x},u_{D,x},u_{A,x})$ es
continua en norma uniforme en $(0,1)$, y para $x,y$ en un compacto con mínimo $x_*>0$ y
$\Lambda=\max\lambda$,
$\|\Theta_x-\Theta_y\|_\infty\le\frac{2C_\Lambda}{x_*^2}|\lambda(1-x)-\lambda(1-y)|+\frac2{x_*^3}|x-y|$.
